\documentclass[12pt, a4paper]{article}

\usepackage{amsmath}
\usepackage{amssymb}
\usepackage{amsthm}
\usepackage{geometry}
\usepackage{titlesec}

\geometry{margin=2.5cm}

% Custom theorem/definition styles
\newtheoremstyle{propertystyle}
  {10pt}{10pt}{\normalfont}{0pt}{\bfseries}{.}{.5em}{}

\theoremstyle{propertystyle}
\newtheorem{property}{Property}
\newtheorem{definition}{Definition}
\newtheorem{theorem}{Theorem}

\titleformat{\section}{\large\bfseries}{}{0em}{}[\titlerule]
\titleformat{\subsection}{\normalsize\bfseries}{}{0em}{}

\title{\textbf{Why Dividing by a Fraction Means Multiplying by Its Reciprocal}}
\author{0xHaru}
\date{\today}

\begin{document}

\maketitle
\thispagestyle{empty}

\begin{property}[Invariant Property]
  Multiplying or dividing both the dividend and the divisor by the same
  non-zero number leaves the quotient unchanged.
\end{property}

\noindent Formally, for all $a, b, k \in \mathbb{Z}$ with $b \neq 0$ and $k \neq 0$:
\[
  \frac{a}{b} = \frac{a \cdot k}{b \cdot k} = \frac{a / k}{b / k}
\]

\begin{definition}[Neutral Element]
  The neutral element of division is $1$.
\end{definition}

\noindent That is, for any $a \in \mathbb{Z}$:
\[
  \frac{a}{1} = a
\]

\begin{theorem}
  Dividing any number by a fraction is equivalent to multiplying that number
  by the reciprocal of the fraction. That is, for all $a, c, d \in \mathbb{Z}$
  with $c \neq 0$ and $d \neq 0$:
  \[
    \dfrac{a}{\;\dfrac{c}{d}\;} = a \cdot \dfrac{d}{c}
  \]
\end{theorem}

\begin{proof}
Let $a, c, d \in \mathbb{Z}$, with $c \neq 0$ and $d \neq 0$, so that
$\dfrac{c}{d}$ is a well-defined, non-zero fraction.

\medskip
\noindent We apply the \textbf{Invariant Property of Division}: multiplying both the
dividend and the divisor by the same non-zero number does not change the quotient.
Here we choose to multiply both by $\dfrac{d}{c}$, the reciprocal of the
divisor $\dfrac{c}{d}$:

\[
  \dfrac{\;a\;}{\;\dfrac{c}{d}\;}
  \;=\;
  \dfrac{\;a \cdot \dfrac{d}{c}\;}
        {\;\dfrac{c}{d} \cdot \dfrac{d}{c}\;}
  \;=\;
  \dfrac{\;a \cdot \dfrac{d}{c}\;}
        {\;1\;}
  \;=\;
  a \cdot \dfrac{d}{c}
\]

\medskip
\noindent The key step is that $\dfrac{c}{d} \cdot \dfrac{d}{c} = \dfrac{cd}{dc} = 1$,
since a fraction multiplied by its reciprocal always equals $1$ (the neutral element
of division). The result follows immediately.
\end{proof}

\newtheorem{corollary}{Corollary}
\begin{corollary}
  Dividing a fraction by a second fraction is equivalent to multiplying the
  first by the reciprocal of the second. For all $a, b, c, d \in \mathbb{Z}$
  with $b \neq 0$, $c \neq 0$, and $d \neq 0$:
  \[
    \dfrac{\;\dfrac{a}{b}\;}{\;\dfrac{c}{d}\;}
    \;=\;
    \dfrac{a}{b} \cdot \dfrac{d}{c}
  \]
\end{corollary}

\begin{proof}
  This follows directly from Theorem~1 by setting $a \leftarrow \dfrac{a}{b}$.
\end{proof}

\end{document}
